\begin{enumerate}

\item El área es $A = \pi d^2/4$.

\item La tensión media es $\sigma = F/A$.

\item La ley de Hooke da $\varepsilon = \sigma/E$ y el alargamiento es $\delta = \varepsilon L$.

\item El coeficiente de seguridad respecto al límite elástico mínimo es $430/\sigma$.

\end{enumerate}

\noindent\textbf{Operaciones y resultados intermedios.}

\begin{enumerate}

\item Área de la sección.

\[\fitmath{y_{1}=\frac{\pi\cdot {0.02}^{2}}{4}\approx 0{,}0003141593\,\mathrm{m^{2}}}\]

\item Tensión normal sigma = F/A.

\[\fitmath{y_{2}=\frac{\frac{60000}{\frac{\pi\cdot {0.02}^{2}}{4}}}{1000000}\approx 190{,}9859\,\mathrm{MPa}}\]

\item Deformación epsilon = sigma/E.

\[\fitmath{y_{3}=\frac{\frac{60000}{\frac{\pi\cdot {0.02}^{2}}{4}}}{7.17\times10^{10}}\approx 0{,}002663681\,}\]

\item Coeficiente de seguridad respecto al límite elástico.

\[\fitmath{y_{4}=\frac{430}{\frac{\frac{60000}{\frac{\pi\cdot {0.02}^{2}}{4}}}{1000000}}\approx 2{,}251475\,}\]

\item Alargamiento total delta = epsilon L.

\[\fitmath{y_{5}=\frac{\frac{60000}{\frac{\pi\cdot {0.02}^{2}}{4}}}{7.17\times10^{10}}\cdot 1{,}5\cdot 1000\approx 3{,}995522\,\mathrm{mm}}\]

\end{enumerate}
